\documentclass[12pt]{article}
\usepackage[margin=1in]{geometry}
\usepackage{enumitem}
\usepackage{titlesec}
\usepackage{parskip}
\usepackage{graphicx}
\usepackage{hyperref}
\usepackage{fontenc}

\title{\textbf{Cox 1997 RG Notes}\\
\large Cox, Gary W. \textit{Making Votes Count: Strategic Coordination in the World's Electoral Systems}.\\
Cambridge: Cambridge University Press, 1997.}
\author{}
\date{}

\begin{document}

\maketitle

\section{Main Idea}



%--------------------------------------------------
\section{General Notes}
%--------------------------------------------------

\begin{itemize}[leftmargin=*, itemsep=4pt]
    \item Contributions: Elections as a coordination problem, empirical tests based on constituency (rather than national aggregate) level.
    \item Elections represent coordination problems: parties (and voters) must coalesce around certain candidates, and do so both by and in light of strategic voting. Yet, this strategic coordination does not always have clear outcomes for winners/losers. Successful coordination often requires an ideological tradeoff that is often unappealing to strictly-ideological factions.
    \item Elections are by and large market affairs, where ``market-clearing'' expectations balance the supply and demand for a number of candidates. Yet, election behavior does not always follow equilibrium outcomes, and in fact a great deal of electoral politics happens outside of these confines or in otherwise unstable institutional settings
    \item Duvergian critiques: broad sociological critiques largely stemming from accusations of oversimplification or ``institutional determinism'': institutions determine this outcome regardless of cultural or social norms, behavior, context, etc.
    \begin{itemize}
        \item Is the direction of causality correct? Could it be the case that party systems determine electoral systems, rather than the inverse? Citations listed on pg. 15
        \item Is the electoral system truly the variable of interest, or might more emphasis fall on the number and structure of social cleavages?
    \end{itemize}
    Quite a bit of this also rises out of questions of endogeneity over electoral systems -- they are, of course, endogenous in initial construction, but what about their long-term forms? Presumably parties can shape electoral institutions over time in accordance with their own strategic visions
    \item Three stages when accounting for the level of vote or seat concentration in some polity:
    \begin{enumerate}
        \item Translation of social cleavages into partisan preferences;
        \item Translation of partisan preferences into votes;
        \item Translation of votes into seats.
    \end{enumerate}
    \item Core construction is the ``$M + 1$ Rule'': if a district offers $M$ seats, we should expect to see consolidation down to $M+1$ viable candidates/parties/etc. participating in that election. Duverger's Law is subsequently a simplification of this stratified between $M = 1, M > 1$ (majoritarianism vs. proportionalism, respectively).
\end{itemize}


\section{Important Specific Factual Claims}

\begin{itemize}[leftmargin=*, itemsep=4pt]
    \item Identifies an \emph{a priori} expectation on the number of parties competing, compared between legislative chambers within a country
    \item Limitation of Duverger: what if cleavages are geographically concentrated? A party can have a stronghold within a small region, but still be relegated to third-party status overall. Think SNP; Duverger's should predict a combination with Labour (but is this potentially stymmied by the possibility of coalition formation?)
\end{itemize}


\section{Connections}

\begin{itemize}
    \item Heavy framing based on the Gibbard-Satterthwaite theorem, but notes that its result -- though quite rigorous -- can be a bit more difficult to apply to political science (just provides a binary Y/N on strategic voting)
    \item Highlights research gap in interaction between electoral systems and social cleavages (though may have been filled in past 30 years)
\end{itemize}


\section{My Critiques}

\begin{itemize}
    \item
\end{itemize}


\end{document}
